\documentclass{article}
\usepackage{amsmath}
\usepackage{fullpage}
\usepackage{amssymb}
\usepackage{booktabs}
\usepackage{array}
\usepackage{tabularx}
\usepackage{hyperref}

\newcolumntype{Y}{>{\raggedright\arraybackslash}X}

\title{Certifying a Source-Based Formalization:\\
the Converse Direction of the Bender--Suzuki Theorem}
\author{}
\date{}

\begin{document}

\maketitle

\section{Certification Without Reading the Proofs}
\label{sec:certification}

This section describes a certification procedure for a completed formalization
and reports its application to the converse direction of the Bender--Suzuki
theorem. The procedure addresses a validation obligation that is distinct from
proof correctness and that is not discharged by the kernel. Section
\ref{sec:validation-obligation} states the obligation. Section
\ref{sec:artifact} describes the challenge/solution artifact that bounds the
material a reader must examine. Section \ref{sec:certificate} presents the
certificate itself, and Section \ref{sec:limits} states the point at which the
certificate ceases to constrain the formalization. Sections
\ref{sec:reading-list} and \ref{sec:not-certified} record, respectively, what a
reader must read and what the procedure does not establish.

\subsection{The Validation Obligation}
\label{sec:validation-obligation}

Every result reported below is machine-checked. On each of the principal
theorems, the axiom census returns exactly
\[
  \{\mathsf{propext},\ \mathsf{Classical.choice},\ \mathsf{Quot.sound}\},
\]
the three axioms of ordinary mathematics in Lean. The development contains no
occurrence of \verb|sorry|, no additional \verb|axiom|, and no invocation of
\verb|native_decide|.

Machine-checking a proof does not, however, certify a \emph{statement}. A
statement may be vacuous, or may quantify over objects other than those
intended, and may nevertheless be proved. Consequently, the obligation
addressed here is not whether the proofs are correct, which the kernel
determines, but rather how much material a reader must examine before accepting
that the statements express what they are claimed to express. We refer to the
material that a reader must examine for this purpose as the \emph{audit
surface}.

This obligation is aggravated by two properties of a large source-based
formalization. First, a statement is expressed in the vocabulary of the
development, so a reader who does not already trust that vocabulary cannot
evaluate the statement in isolation. Second, the vocabulary is extensive: the
development considered here defines several thousand declarations, of which the
statements of interest use only a small number, and identifying that number is
not immediate from inspection of the source files.

\subsection{The Challenge/Solution Artifact}
\label{sec:artifact}

The artifact is a challenge/solution pair in the sense of the \textsc{Comparator}
tool of the Lean FRO.\footnote{\url{https://github.com/leanprover/comparator}}
It consists of three Lean modules and a configuration file, summarized in
Table \ref{tab:files}.

\begin{table}[t]
\centering
\begin{tabularx}{\textwidth}{l r Y}
\toprule
file & lines & role \\
\midrule
\texttt{comparator/Defs.lean} & 253 & the definitions, stated over Mathlib alone \\
\texttt{comparator/Challenge.lean} & 56 & four statements, each with proof \texttt{sorry} \\
\texttt{comparator/Solution.lean} & 89 & the same four statements, proved \\
\texttt{comparator/config.json} & --- & theorem names and permitted axioms \\
\bottomrule
\end{tabularx}
\caption{The challenge/solution artifact. The first two files constitute the
audit surface.}
\label{tab:files}
\end{table}

\paragraph{Independence from the development.}
The challenge does not import the development it certifies. Its subject matter,
namely strong embedding, Hypothesis (A), and the three groups appearing in the
conclusion of the classification, is defined in \verb|Defs.lean| from Mathlib
primitives alone. Let the \emph{statement cone} of a set of declarations denote
the transitive closure of the constants occurring in their types, closed under
types, definition values, and inductive constructors. This is the set of
declarations that \textsc{Comparator} compares, and therefore the set whose
meaning a reader must accept. For the four challenge statements the cone
contains 3786 constants, of which 33 are declared in \verb|Defs.lean|, and 27
after auxiliary proof obligations of the form \verb|_proof_*| are discarded.
\emph{No constant of the development occurs in the cone.}

An earlier version of the artifact named the development's own
\verb|HypothesisA| and \verb|suzukiConclusion|. That version carried 37 imported
definitions that a reader was required to accept without inspection. It is
retained in the version-control history.

\paragraph{Shared rather than duplicated definitions.}
The usual arrangement requires the solution to restate the challenge's
definitions verbatim, and the comparison procedure then verifies that the two
copies agree. This arrangement does not scale to the present setting. The two
files have different import sets, since the solution imports the entire
development, and instance resolution in Lean is sensitive to the import set.
Three divergences were observed on the duplicated version:
\begin{enumerate}
  \item[\textsc{[E1]}] the element \(\bot\) occurring in \texttt{Disjoint}
  elaborated through \texttt{BoundedOrder.toOrderBot} in the challenge and
  through \texttt{CompletePartialOrder.toOrderBot} in the solution;

  \item[\textsc{[E2]}] the expression \verb|GaloisField 2 m| resolved the
  instance \(\mathsf{Fact}(\mathsf{Nat.Prime}\ 2)\) to \verb|Nat.fact_prime_two|
  in one file and to the development's own instance in the other;

  \item[\textsc{[E3]}] the finiteness instance for the linear group, left
  implicit in the statement, resolved in the challenge and in the solution
  respectively through
  \[
    \texttt{Subgroup.finiteIndex\_of\_finite}
    \qquad\text{and}\qquad
    \texttt{Subgroup.finiteIndex\_center} .
  \]
\end{enumerate}
The comparison is syntactic, so each of these divergences would have been
rejected. Byte-identical source therefore does not yield byte-identical
elaboration. Divergences \textsc{[E1]} and \textsc{[E2]} are eliminated by
placing the definitions in a module imported by both files, so that the two
occurrences are the same constant rather than two copies required to agree by
accident. Divergence \textsc{[E3]} is eliminated by binding the instance
explicitly in the statement.

This arrangement does not weaken the artifact. The module \verb|Defs.lean| lies
in the import closure of the challenge, hence within the audit surface, and the
solution cannot vary it.

\paragraph{Admissibility of imports on the solution side.}
The solution imports the entire development. This is admissible by construction
rather than by convention, for the following reasons.
\begin{enumerate}
  \item The export step emits the transitive closure of the named declarations.
  \item The replay step constructs a fresh empty environment and replays that
  closure through the kernel. Nothing is trusted from precompiled artifacts;
  every constant in the proof cone is typechecked again from the exported term.
  \item The axiom check enforces the permitted set across the same closure, so
  an occurrence of \verb|sorry| anywhere in the imported development would
  appear as \(\mathsf{sorryAx}\) and be rejected.
  \item For the four targets, only the name, the universe parameters, and the
  type are compared. The proof term is not compared and may therefore originate
  anywhere.
\end{enumerate}
In particular, nothing requires the solution to use the development. This
omission is deliberate and does not constitute a gap, since the solution
\emph{is} the development, so whatever the solution discharges, the development
proves.

\paragraph{The trusted base.}
Under the first assumption of the comparison tool, the trusted set is the import
closure of the challenge together with the build configuration. The audit
surface therefore consists of \verb|Defs.lean| (253 lines),
\verb|Challenge.lean| (56 lines), and the twelve-line addition to
\verb|lakefile.toml|, and of nothing further. Mathlib is assumed; no part of the
development is.

The configuration declares no definition holes. A definition hole is
exploitable: a solution could define strong embedding as the constantly true
predicate and discharge every obligation trivially. Since no definition is
declared as a hole, every definition is reached by the comparison procedure as
part of the statement cone and must agree with the solution's occurrence in full.

\paragraph{The four obligations.}
The challenge states the four obligations \textsc{[V1]}--\textsc{[V4]} listed
below.
\begin{enumerate}
  \item[\textsc{[V1]}] Hypothesis (A) implies that the point stabilizer is
  strongly embedded.

  \item[\textsc{[V2]}] Hypothesis (A) is satisfiable. Specifically,
  \(\mathrm{PSL}_2(2^k)\) satisfies it for every \(k\ge 2\). Prior to this work
  no declaration in the development ever constructed an instance of Hypothesis
  (A); every occurrence consumed one, so the classification theorem might have
  been vacuously true.

  \item[\textsc{[V3]}] The Suzuki branch is attained: \(\mathrm{Sz}(2^{2m+1})\)
  satisfies Hypothesis (A) for every \(m\ge 1\).

  \item[\textsc{[V4]}] The unitary branch is attained: for every \(k\ge2\) there
  exist a field of order \((2^k)^2\) whose fixed field has order \(2^k\), and a
  Hermitian form with antidiagonal Gram matrix, such that the associated
  \(\mathrm{PSU}_3\) satisfies Hypothesis (A).
\end{enumerate}
Composing \textsc{[V1]} with each of \textsc{[V2]}--\textsc{[V4]} establishes
that each of the three groups possesses a strongly embedded subgroup. Every
object supplied by the solution is existentially quantified in the challenge,
namely the set acted upon, the action, and the subgroups \(H\), \(D\), \(Q\)
together with the involution \(t\). Consequently the solution cannot substitute
a group of its own choosing.

\paragraph{Recognizability of the definitions.}
The definition of \(\mathrm{PSL}_2(2^k)\) occupies one line, namely
\[
  \texttt{Matrix.ProjectiveSpecialLinearGroup (Fin 2) (GaloisField 2 k)} .
\]
The unitary group is the isometry group of a Hermitian form, restricted to
determinant one and mapped into the projective general linear group. Both are
recognizable on inspection.

The Suzuki group is not. Its definition is an explicit generating set of
\(4\times4\) matrices involving a Tits exponent \(2^{m+1}\), and no reader
should be required to recognize \(\mathrm{Sz}(q)\) in it. No reader is so
required. The obligations state that the group satisfies Hypothesis (A) and
therefore possesses a strongly embedded subgroup; Bender's classification then
places it within the three families, and the group orders separate it from the
remaining two. An unrecognizable definition thus certifies itself.

We observe further that proof terms occurring \emph{inside} definitions, such as
the closure obligations of the unitary subgroup and the determinant computations
of the Suzuki generators, carry no mathematical content: a subgroup is
determined by its carrier. The audit cost of a sixty-line definition with a
four-line carrier is four lines.

\paragraph{Verification performed.}
The following checks were carried out.
\begin{enumerate}
  \item The module \verb|Defs.lean| elaborates against Mathlib alone.
  \item The module \verb|Challenge.lean| elaborates with exactly four
  occurrences of \verb|sorry| and no errors.
  \item The module \verb|Solution.lean| elaborates with no occurrence of
  \verb|sorry|.
  \item Printing the four statement types with \verb|pp.all| from each module
  and comparing the results yields 9305 identical lines, which is the property
  that the comparison tool itself enforces.
  \item The command \verb|#print axioms| on each of the four solution theorems
  returns the permitted set.
  \item The statement cone was enumerated by a traversal over types, definition
  values, and inductive constructors, and contains no constant of the
  development.
\end{enumerate}
The comparison tool was executed and reports that the solution is accepted. It
builds and exports each module within the sandbox, compares the four statements,
checks the axioms, and replays the solution through a fresh kernel.

A successful run establishes nothing unless the procedure is capable of
rejection. Two negative controls were therefore performed. Replacing one
solution proof by \verb|sorry| yields the diagnostic \emph{illegal axiom
detected}; weakening one hypothesis yields \emph{challenge and solution theorem
statement do not match}. Both are rejections of the kind on which the artifact
relies.

Four adaptations of the environment were required, and we record them since each
is a deviation from the intended configuration. First, the distributed exporter
targets a later version of Lean than the development and cannot read its
compiled artifacts; it was rebuilt from the same source against the version in
use. Second, the current release of the sandboxing tool requires a Landlock
interface later than the one this kernel provides, so an earlier release was
used. Third, the tool grants read access to the system directories but not
execute access, and the automatic detection of library dependencies in that
release omits the dynamic loader; a wrapper supplies execute rights on the
system library directories, leaving the write restrictions unchanged. That the
restrictions remain in force was verified separately, by confirming that reads
and writes outside the permitted paths are refused. Fourth, the build system's
artifact cache lies outside the writable set, so the home directory was directed
to a location within it.

The tool further requires that the solution not have been compiled in the
checking environment beforehand, since a challenge compiled alongside an
adversarial solution might otherwise pass unnoticed. This requirement concerns
the environment of the party performing the check rather than the artifact, and
is discharged by obtaining the artifact and running the comparison in a fresh
checkout. It was also discharged here, by deleting the compiled artifacts of the
three modules and repeating the comparison, so that the tool rebuilt them from
source within the sandbox; its own order builds and exports the challenge before
the solution is compiled. The outcome was unchanged.

\subsection{Strong Embedding as a Certificate}
\label{sec:certificate}

A subgroup \(M\) of a finite group \(G\) is \emph{strongly embedded} if
\[
  M\neq G,
  \qquad
  2\mid\lvert M\rvert,
  \qquad\text{and}\qquad
  2\nmid\lvert M\cap gMg^{-1}\rvert
  \quad\text{for every }g\in G\setminus M .
\]
The Lean declaration is a structure carrying exactly these three fields, named
\verb|ne_top|, \verb|even_card|, and \verb|odd_inter|, and occupies four lines.

This notion is absent from Mathlib and therefore constitutes new code that must
be read. It is, however, expressed entirely in Mathlib vocabulary: each of the
six symbols occurring in it belongs to Mathlib, and its import list consists of
four Mathlib modules with no dependency on the development, a property that the
build verifies. Conjugation is written as the image of \(M\) under Mathlib's
conjugation automorphism, which denotes \(gMg^{-1}\); this is verified in two
ways, definitionally against the pointwise action, and by the membership
criterion \(x\in gMg^{-1}\iff g^{-1}xg\in M\).

We now state the certification argument. It is proved that each of the three
families possesses a strongly embedded subgroup, and, more generally, that the
normal subgroup produced by the classification always does, with no faithfulness
hypothesis required.

A group possessing a strongly embedded subgroup \emph{and} of \(2\)-rank at
least two is severely constrained. By Bender's theorem, which is external to the
development, such a group \(G\) possesses a normal subgroup \(L\) of odd index
with \(L/O(L)\) isomorphic to one of \(\mathrm{PSL}_2(2^n)\),
\(\mathrm{Sz}(2^{2n+1})\), or \(\mathrm{PSU}_3(2^n)\). The rank condition is not
decoration: without it the conclusion is false, since in any group whose Sylow
\(2\)-subgroups are cyclic or generalized quaternion, the normalizer of a Sylow
\(2\)-subgroup is strongly embedded. Both hypotheses are available here, since
the rank condition is a field of Hypothesis (A) and is therefore established for
each of the three families alongside the strong embedding.

A reader who accepts that classical theorem, and who has read the four lines
displayed above, thereby obtains the following conclusion: whatever the three
matrix-group definitions of the development actually denote, they denote groups
in Bender's list. They cannot be trivial, cannot be solvable, and cannot be
matrix groups of the wrong shape. The argument therefore constitutes an
inexpensive check on several thousand lines of matrix-group construction,
obtained at the cost of reading a four-line structure.

\subsection{Limits of the Certificate}
\label{sec:limits}

The certificate establishes \emph{membership} in Bender's list and not
\emph{which member} is obtained. A reader is entitled to know precisely where
the argument ceases to constrain the formalization.

Suppose that the Suzuki group had been defined erroneously as
\(\mathrm{PSL}_2(2^{2m+1})\). It would still possess a strongly embedded
subgroup and would still satisfy Hypothesis (A), and nothing in Section
\ref{sec:certificate} would detect the error. The coincidence is not
hypothetical, as Table \ref{tab:degree65} indicates.

\begin{table}[t]
\centering
\begin{tabular}{lrr}
\toprule
group & degree & order \\
\midrule
\(\mathrm{Sz}(8)\)        & 65 & 29\,120 \\
\(\mathrm{PSL}_2(64)\)    & 65 & 262\,080 \\
\(\mathrm{PSU}_3(4)\)     & 65 & 62\,400 \\
\bottomrule
\end{tabular}
\caption{Three groups that are doubly transitive of the same degree and each of
which satisfies Hypothesis (A).}
\label{tab:degree65}
\end{table}

All three groups are doubly transitive of degree \(65\) and each satisfies
Hypothesis (A). Degree alone, and therefore Hypothesis (A) alone, does not
distinguish them; the group order does. Coverage of the orders is uneven, and we
record the residual gap explicitly.
\begin{itemize}
  \item For the Suzuki family, the order
  \(\lvert\mathrm{Sz}(q)\rvert=(q^2+1)q^2(q-1)\) is a stated theorem of the
  development. The family is therefore fully determined.
  \item For the unitary family, the quantities \(\lvert\Omega\rvert=q^3+1\) and
  \(\lvert U\rvert=q^3(q^2-1)/\gcd(q+1,3)\) are stated. The group order follows
  by the orbit-stabilizer theorem but is not itself a named theorem.
  \item For the linear family, the quantities \(\lvert Q\rvert=q\),
  \(\lvert D\rvert=q-1\), and \(\lvert\Omega\rvert=q+1\) are stated. The order
  again follows but is not named.
\end{itemize}
Supplying the two missing order statements would close this gap. They are
corollaries of facts already established and do not constitute new mathematics.

\subsection{Reading List}
\label{sec:reading-list}

A reader who wishes to accept the principal claims, and who is prepared to trust
the Lean kernel for everything else, must read the statements listed in Table
\ref{tab:reading}, and no proofs.

\begin{table}[t]
\centering
\begin{tabularx}{\textwidth}{Y l r}
\toprule
content & location & size \\
\midrule
strong embedding, Hypothesis (A), the three groups &
\texttt{comparator/Defs.lean} & 253 lines \\
the four obligations & \texttt{comparator/Challenge.lean} & 56 lines \\
\bottomrule
\end{tabularx}
\caption{The audit surface. Its import closure is Mathlib and nothing further.}
\label{tab:reading}
\end{table}

The total is 309 lines, whose import closure is Mathlib and nothing further. No
file of the development appears in the list, which is the purpose of stating the
challenge over Mathlib. Of the 253 lines of definitions, approximately 60
discharge subgroup-closure and determinant obligations occurring inside
definitions and carry no mathematical content; the substantive reading is
therefore nearer 190 lines. The remaining new material, approximately 2700 lines,
consists of proofs and is the concern of the kernel rather than of the reader.

\subsection{Claims Not Supported}
\label{sec:not-certified}

For completeness, we record the claims that the preceding argument does
\emph{not} support.
\begin{enumerate}
  \item That Bender's theorem holds. It is used in Section
  \ref{sec:certificate} as external knowledge and is not formalized in the
  development.

  \item That Hypothesis (A) is equivalent to the existence of a strongly
  embedded subgroup. Only the forward implication is established; the converse
  is the substantive direction of Bender's paper.

  \item That the three families are pairwise non-isomorphic, or that the
  definitions select the intended member of the list. See Section
  \ref{sec:limits}.

  \item That the ambient group of an arbitrary instance of the conclusion
  possesses a strongly embedded subgroup in the absence of the faithfulness
  hypothesis. The statement is presumably true, since for \(G=L\times C\) with
  \(C\) of odd order one may take \(H\times C\), but it is not proved.

  \item Any assertion concerning the classification theorem \emph{as a named
  declaration}. The challenge does not mention it. What the artifact certifies is
  that the development proves the four obligations of Section \ref{sec:artifact}.
  The equivalence between Hypothesis (A) and the conclusion of the
  classification is established in the development but lies outside the audit
  surface.
\end{enumerate}

\section{Relation to the Sources}
\label{sec:sources}

This section is a source-critical account of the formalization. Section
\ref{sec:provenance} identifies the results formalized and their provenance.
Section \ref{sec:agreement} records the respects in which the formalization
follows its sources closely, and Section \ref{sec:departures} classifies the
respects in which it departs from them.

\subsection{Scope and Provenance}
\label{sec:provenance}

The pre-existing theorem of the development is Peterfalvi, \emph{Character
Theory for the Odd Order Theorem}, Part II, Theorem A: the Bender--Suzuki
classification in its Zassenhaus form, that is, a classification of the doubly
transitive groups of even degree whose point stabilizer splits as a normal
subgroup of even order by a complement of odd order. Mathematically this is
Suzuki's 1962 determination of that class, together with Bender's subsequent
simplification. Its classical inputs are themselves formalized rather than
assumed: Huppert, \emph{Endliche Gruppen I}, II.10.12 for the unitary family,
and Huppert--Blackburn XI.3.1 and XI.3.3 for the Suzuki family. The linear
family rests on Mathlib.

The work reported here contributes three items, none of which is a transcription
of a numbered statement in any source.
\begin{enumerate}
  \item \emph{Witnesses}. Each of the three families appearing in the conclusion
  is shown to satisfy the hypothesis.
  \item \emph{The converse}, and hence an equivalence between the hypothesis and
  the conclusion.
  \item \emph{Strong embedding}, together with the implication that the point
  stabilizer under Hypothesis (A) is strongly embedded.
\end{enumerate}

\subsection{Agreement with the Sources}
\label{sec:agreement}

\paragraph{Field-for-field transcription of the hypothesis.}
Hypothesis (A1) is transcribed with twelve fields, in the order and shape of
Peterfalvi's Part II hypothesis: double transitivity, the identification of
\(H\) as a point stabilizer, the assertion that \(t\) is an involution outside
\(H\), the equation \(D=H\cap H^t\), the semidirect decomposition \(H=Q\rtimes D\)
expressed as five separate conditions, and finally the parity conditions on
\(\lvert Q\rvert\) and \(\lvert D\rvert\). The remaining two fields are
faithfulness and the condition that the \(2\)-rank is at least two.

\paragraph{Preservation of the conjugation convention.}
The development defines \(H^g=g^{-1}Hg\), following the exponent convention of
the text rather than Mathlib's opposite convention. The discrepancy is recorded
in the documentation string. The convention occurs in the statement of the
theorem and not merely in its proof.

\paragraph{Explicit equivariance in the conclusion.}
The conclusion asserts more than an abstract isomorphism onto a member of the
list. Each of the three action models carries a group isomorphism, a permutation
representation pinned to the linear action, and a bijection intertwining them.
The three sets acted upon are the textbook ones: the projective line, the Tits
ovoid in \(\mathrm{PG}(3,q)\) in the coordinates of Huppert--Blackburn, and the
isotropic points of a Hermitian form with antidiagonal Gram matrix.

\subsection{Departures from the Sources}
\label{sec:departures}

We classify the departures into the seven cases D1--D7.

\paragraph{D1: strong embedding is not the hypothesis assumed.}
The Bender--Suzuki theorem is frequently quoted in the form ``if \(G\) possesses
a strongly embedded subgroup, then \ldots''. This is not the hypothesis that the
development assumes, and the term did not occur anywhere in the development
prior to this work. Hypothesis (A) assumes the full Zassenhaus configuration,
which is strictly stronger than strong embedding. Recovering the configuration
from a strongly embedded subgroup alone is the substantive content of Bender's
1971 paper and is not formalized. Only the forward implication is established.

\paragraph{D2: the converse has no source statement.}
Peterfalvi states Theorem A in one direction only. That the three families
satisfy Hypothesis (A) is classical and universally assumed, but is not a
numbered result available for transcription; it is assembled from the classical
inputs. The passage from a normal subgroup of odd index to the ambient group is
likewise standard but carries no number in any source, and is proved from first
principles.

\paragraph{D3: a hypothesis absent from the sources.}
The equivalence requires faithfulness of the action on the classification side.
This requirement is not stylistic. Without it the statement is false: for any
group \(C\) of odd order acting trivially, the product \(L\times C\) satisfies
the conclusion of the classification but fails faithfulness. Textbook statements
do not exhibit this phenomenon because they do not attempt the converse.

\paragraph{D4: a convention that departs deliberately.}
The definition of strong embedding uses Mathlib's conjugation convention and not
the development's. The justification is that this is a general-purpose
definition rather than a transcription of a numbered statement, so the argument
from fidelity to the source does not apply, and legibility against Mathlib is of
greater value. The two conventions are interchangeable in this instance, since
the conjugating element ranges over the complement of the subgroup and that
complement is closed under inversion. A bridging lemma records the equivalence.

\paragraph{D5: proof strategies that depart from the sources.}
In the Suzuki case, Huppert--Blackburn establish the complement conditions by
computing the action of root elements on the ovoid. We instead derive
\(\lvert Q\rvert=q^2\) from the root-closure description and
\(\lvert D\rvert=q-1\) from the orbit-stabilizer theorem together with double
transitivity. The intersection condition then follows from the coprimality of
\(q^2\) and \(q-1\), and the join condition from a cardinality count. The ovoid
computation is avoided entirely. The mathematics is unchanged; the route is not
that of the source.

In the lifting argument, normality of \(Q\) in the larger stabilizer is obtained
by showing that \(Q\) absorbs every \(2\)-element of the relevant intersection,
by way of the quotient by \(Q\), whose order is odd. The customary route invokes
Sylow uniqueness.

\paragraph{D6: departures forced by formalization.}
The set acted upon is an independent type carrying an action, rather than a coset
space. The unitary branch of the conclusion quantifies existentially over the
field and the form rather than naming a specific field. This existential form is
weaker than the textbook statement and materially complicated the converse: a
witness for one particular form does not suffice, so the unitary development was
generalized to an arbitrary triple of the stated shape. Finally, the conclusion
supplies explicit data where the sources assert that the group ``is isomorphic
to, acting naturally''. This is a strengthening, and it is what renders the
converse provable.

\paragraph{D7: a second transcription in Mathlib vocabulary.}
For the certification described in Section \ref{sec:certification}, Hypothesis
(A) and the three groups are transcribed a second time, in the audit surface,
using Mathlib primitives alone. The copy is field for field, and the solution is
required to discharge the development's version from it, so that a weakened copy
would fail rather than pass undetected. The transcription forces two further
departures. First, the conjugation operator is written directly in terms of
Mathlib's conjugation automorphism, since the development's abbreviation for it
is not a Mathlib declaration. Second, the three groups are transcribed under new
names; each is definitionally equal to the development's version, and this is
verified by reflexivity.

\end{document}
